\section{The case that the sample size is fixed}
\label{sec:fixed}

Theorem~\ref{thm:5} requires $n_i\to\infty$. We show that this is not a technical
restriction but is essential.

\begin{proposition}
\label{prop:3}
Assume \textup{(E)} with $c\in(0,1)$ fixed and let $N$ be fixed. Consider any
procedure which uses a projection $\hat P$ of fixed rank estimated from the
training data. Then the angle between $\hat h$ and the true spiked direction $h$
does not tend to zero, that is, $|\hat h^Th|\to\varrho<1$ in probability. Hence
the residual spiked component of the new observation after the projection
satisfies
\[
\bigl\|(I-\hat h\hat h^T)\sqrt{\lambda}\,z_0 h\bigr\|^2
=\lambda z_0^2(1-\varrho^2)+o_P(\theta)=O_P(\theta),
\]
so that the spike cannot be removed by the projection and the projection-based
SC-SVM does not hold the consistency property.
\end{proposition}

\begin{proof}[Outline of proof]
The region $\lambda/\tr(\Sigma_*)\to c/(1-c)\in(0,\infty)$ corresponds to the
subspace inconsistency region in the HDLSS PCA theory of \citet{jung2009}, in
which the sample principal component direction is not consistent and the angle has
a non-degenerate limit. Intuitively, the information about the spiked direction
contained in the sample consists of the $N$ coefficients $\{z_{ijs}\}$ only, on
which the isotropic noise of order $\sqrt{\tr(\Sigma_*)}$ is superposed, so that
the signal-to-noise ratio remains a constant when $N$ is fixed. Since the
estimation error of $\hat h$ remains at a constant proportion, the spiked
component $\sqrt{\lambda}\,z_0 h$ of $x_0$, whose magnitude is of order
$\sqrt{\theta}$, also remains at a constant proportion after the projection. Then
the argument of Theorem~\ref{thm:3} is applicable and the misclassification rate
does not tend to zero.
\end{proof}

\begin{remark}
\label{rem:7}
Proposition~\ref{prop:3} is a statement about the class of projection-based
procedures. We have not proved a minimax type lower bound which states that any
classification rule based on the training data does not hold the consistency
property when $N$ is fixed and $\delta<\infty$. We expect that it is reduced to the
positivity of the Bayes risk in the limiting experiment given in
Theorem~\ref{thm:1}, but it remains open. On the other hand, when $\delta\to\infty$,
we can see from~\eqref{eq:7.2} that the consistency property is recovered even for
fixed $N$.
\end{remark}
